\documentclass[11pt, a4paper]{article}

%% ── Packages ──────────────────────────────────────────────────────────────────
\usepackage[utf8]{inputenc}
\usepackage[T1]{fontenc}
% \usepackage{lmodern}  % disabled: font not fully installed

% Mathematics
\usepackage{amsmath, amssymb, amsthm, mathtools}
\usepackage{bm}              % bold math symbols

% Figures & graphics
\usepackage{graphicx}
\usepackage{tikz}
\usepackage{pgfplots}
\pgfplotsset{compat=newest}

% Layout & formatting
\usepackage[margin=2.5cm]{geometry}
\usepackage{parskip}         % paragraphs separated by vertical space
\usepackage{enumitem}
\usepackage{booktabs}        % professional tables
% \usepackage{microtype}       % disabled: needs scalable fonts

% References & links
\usepackage[colorlinks=true, linkcolor=blue!60!black, citecolor=green!50!black,
            urlcolor=blue!70!black]{hyperref}
\usepackage{cleveref}        % smart cross-references (\cref)

% Algorithms
\usepackage{algorithm}
\usepackage{algpseudocode}

% Code listings (for referencing C++ implementation)
\usepackage{listings}
\lstset{
  basicstyle=\ttfamily\footnotesize,
  keywordstyle=\bfseries\color{blue!70!black},
  commentstyle=\itshape\color{gray},
  frame=single,
  breaklines=true,
}

%% ── Theorem environments ─────────────────────────────────────────────────────
\theoremstyle{definition}
\newtheorem{definition}{Definition}[section]
\newtheorem{assumption}{Assumption}[section]

\theoremstyle{plain}
\newtheorem{theorem}{Theorem}[section]
\newtheorem{lemma}[theorem]{Lemma}
\newtheorem{proposition}[theorem]{Proposition}
\newtheorem{corollary}[theorem]{Corollary}

\theoremstyle{remark}
\newtheorem{remark}{Remark}[section]

%% ── cleveref names ───────────────────────────────────────────────────────────
\crefname{assumption}{Assumption}{Assumptions}
\crefname{remark}{Remark}{Remarks}
\crefname{definition}{Definition}{Definitions}

%% ── Notation shortcuts ───────────────────────────────────────────────────────
\newcommand{\R}{\mathbb{R}}
\newcommand{\Z}{\mathbb{Z}}
\newcommand{\E}{\mathbb{E}}
\newcommand{\Var}{\operatorname{Var}}
\newcommand{\dd}{\,\mathrm{d}}                   % differential
\newcommand{\ts}[1]{t_{#1}}                       % timestamp shorthand
\newcommand{\niq}{n_{\text{iq}}}                   % queue length (samples)
\newcommand{\freq}{\hat{f}}                        % estimated frequency ratio
\newcommand{\drift}{\hat{\delta}}                  % estimated clock drift

%% ── Title block ──────────────────────────────────────────────────────────────
\title{%
 Latency Controllers\\
  for Real-Time Audio Streaming over IP%
}
\author{%
  % Add author names here
}
\date{\today}

%% ═════════════════════════════════════════════════════════════════════════════
\begin{document}
\maketitle

\begin{abstract}
  % Brief summary:
  % - Problem: maintaining bounded end-to-end latency between independent
  %   audio clocks over a jittery network.
  % - Approach: model-based controller that decouples clock-drift estimation
  %   from jitter-buffer regulation.
  % - Key results / guarantees.
  Vibe-mathed mathematical setting for latency control: with optimality criteria
  (latency bias, latency variance, warp factor smoothness), an SDE model of
  jitter, packet arrivals and sample queue depth, and derived optimal controllers.
\end{abstract}

\tableofcontents
\newpage

%% ══════════════════════════════════════════════════════════════════════════════
\section{Introduction}\label{sec:intro}

% Context: real-time audio streaming (ROC toolkit), sender & receiver with
% independent clocks, network jitter, need for low & stable latency.

% Problem statement: the receiver must continuously adjust its playback
% resampling ratio so that the jitter buffer stays close to a target level,
% despite (a) clock drift, (b) network jitter, and (c) occasional
% discontinuities (session restarts, NTP corrections).

% Contributions of this document.

TODO

%% ══════════════════════════════════════════════════════════════════════════════
\section{System Model}\label{sec:model}

We consider a sender and a receiver, each equipped with an independent
free-running clock.  Neither clock is assumed to be ``correct''; instead
we introduce an absolute reference time~$t \in \R$ that is external to
both devices (one may think of it as an ideal UTC clock that neither
party has access to).

%% ──────────────────────────────────────────────────────────────────────────────
\subsection{Clocks and Drift}\label{sec:clocks}

Let $\tau_A(t)$ denote the reading of the \emph{sender clock} (clock~A)
at absolute time~$t$, and $\tau_B(t)$ the reading of the
\emph{receiver clock} (clock~B).
Each clock runs at a rate that deviates from the reference by a small,
slowly-varying fractional offset:
\begin{align}
  \frac{\dd \tau_A}{\dd t} &= 1 + \alpha(t), \label{eq:clock-a} \\[4pt]
  \frac{\dd \tau_B}{\dd t} &= 1 + \beta(t),  \label{eq:clock-b}
\end{align}
where $\alpha(t)$ and $\beta(t)$ are the instantaneous fractional
frequency offsets of clocks~A and~B respectively (typically on the order
of $10^{-6}$--$10^{-5}$, i.e.\ a few ppm for quartz oscillators).

\begin{assumption}[Nearly-stationary drift]\label{as:drift}
  The offsets $\alpha$ and~$\beta$ are approximately constant over the
  controller's timescale (seconds to minutes) but may drift slowly over
  longer horizons (tens of minutes to hours) due to temperature or
  voltage changes.  Formally, we treat each as a piecewise-linear
  function of~$t$ whose slope is small enough to be neglected within any
  single control epoch.
\end{assumption}

Under \cref{as:drift}, integrating \eqref{eq:clock-a} gives
\begin{equation}\label{eq:tau-a}
  \tau_A(t) \;=\; t + \int_0^{t} \alpha(s)\,\dd s
            \;=\; (1 + \alpha)\,t,
\end{equation}
where the equality holds when $\alpha$ is treated as constant.
An analogous expression holds for $\tau_B$.

The quantity that ultimately matters for clock synchronisation is the
\emph{relative drift} between sender and receiver:
\begin{equation}\label{eq:rel-drift}
  \delta \;=\; \frac{(1+\alpha)-(1+\beta)}{1+\beta}
         \;=\; \frac{\alpha - \beta}{1+\beta}.
\end{equation}

\begin{remark}
  The receiver never observes $\alpha$ or $\beta$ individually---only
  their difference~$\delta$ is identifiable from packet timing.  The
  reference time~$t$ is a modelling convenience; all observable
  quantities will be expressed in receiver-clock time~$\tau_B$.
\end{remark}

%% ──────────────────────────────────────────────────────────────────────────────
\subsection{Packet Emission}\label{sec:emission}

The sender packetises audio into frames of $N$~samples at nominal sample
rate $F_s$ (Hz).  Packet~$n$ ($n = 0, 1, 2, \dots$) is emitted when the
sender clock reads
\begin{equation}\label{eq:send-clock}
  \tau_A^{(\text{send})}(n) \;=\; n\,\Delta, \qquad
  \Delta \;=\; \frac{N}{F_s},
\end{equation}
where $\Delta$ is the \emph{packet period} (seconds of sender-clock time).

The corresponding absolute time of emission is obtained by inverting
\eqref{eq:tau-a}:
\begin{equation}\label{eq:send-abs}
  t^{(\text{send})}(n)
    \;=\; \frac{n\,\Delta}{1 + \alpha}.
\end{equation}

%% ──────────────────────────────────────────────────────────────────────────────
\subsection{Network Transit and Arrival Times}\label{sec:arrival}

Packet~$n$ traverses the network and arrives at the receiver at absolute
time
\begin{equation}\label{eq:arrive-abs}
  t^{(\text{arr})}(n) \;=\; t^{(\text{send})}(n) + d(n),
\end{equation}
where $d(n)$ is the \emph{one-way network delay} for packet~$n$.

We assume that FEC fully recovers all losses, so every emitted packet
eventually arrives.  The delay is decomposed as
\begin{equation}\label{eq:delay-decomp}
  d(n) \;=\; d_0 + \eta(n),
\end{equation}
where $d_0 > 0$ is a baseline (minimum or mean) propagation delay
and $\eta(n)$ is the \emph{jitter}---the stochastic, packet-varying
component of delay.

\begin{assumption}[Jitter model]\label{as:jitter}
  The jitter process $\eta(n)$ may exhibit short-range correlations
  (e.g.\ bursty congestion, queuing effects), but is
  \emph{long-range uncorrelated} in the following sense:
  there exists a mixing horizon $M$ such that for $|i - j| > M$,
  $\eta(i)$ and $\eta(j)$ are approximately independent.

  On a per-packet timescale we define the jitter increments
  \begin{equation}\label{eq:jitter-inc}
    \xi(n) \;=\; \eta(n) - \eta(n-1),
  \end{equation}
  and require only that they are zero-mean with finite variance:
  \begin{equation}\label{eq:jitter-moments}
    \E[\xi(n)] = 0, \qquad
    \Var[\xi(n)] = \sigma_\xi^2 < \infty.
  \end{equation}
  No further distributional assumptions are imposed.
\end{assumption}

\begin{remark}[Packet reordering]
  This model admits packet reordering: if $\xi(n) < -\Delta/(1+\alpha)$,
  then $t^{(\text{arr})}(n) < t^{(\text{arr})}(n-1)$, i.e.\ packet~$n$
  arrives before packet~$n{-}1$.  The jitter buffer at the receiver
  is responsible for restoring packet order before playback.
\end{remark}


\paragraph{Arrival time in receiver clock.}
The receiver observes arrivals in its own clock~$\tau_B$.
Applying \eqref{eq:clock-b} to \eqref{eq:arrive-abs}:
\begin{equation}\label{eq:arrive-recv}
  \tau_B^{(\text{arr})}(n)
    \;=\; (1+\beta)\,t^{(\text{arr})}(n)
    \;=\; (1+\beta)\!\left[\frac{n\,\Delta}{1+\alpha} + d_0 + \eta(n)\right].
\end{equation}

For the \emph{inter-arrival time} observed by the receiver:
\begin{equation}\label{eq:inter-arrival}
  \tau_B^{(\text{arr})}(n) - \tau_B^{(\text{arr})}(n-1)
    \;=\; (1+\beta)\!\left[\frac{\Delta}{1+\alpha} + \xi(n)\right]
    \;=\; \frac{\Delta}{1+\delta} \;+\; (1+\beta)\,\xi(n),
\end{equation}
where the second equality uses $\frac{1+\beta}{1+\alpha} = \frac{1}{1+\delta}$
from \eqref{eq:rel-drift}.

This is the fundamental measurement equation: each inter-arrival
interval reveals the packet period~$\Delta$ scaled by the relative drift
$1/(1+\delta)$, plus a jitter noise term scaled by $(1+\beta)$.

%% ══════════════════════════════════════════════════════════════════════════════
\section{Controller Model}\label{sec:controller-model}

All quantities in this section are expressed in receiver-clock
time, which we denote~$\tau \equiv \tau_B$ for brevity.

%% ──────────────────────────────────────────────────────────────────────────────
\subsection{Arrival Process}\label{sec:arrival-process}

The \emph{cumulative arrival process} is
\begin{equation}\label{eq:arrival-process}
  A(\tau) \;=\; N \cdot \#\bigl\{\,n \ge 0 \;:\;
    \tau_B^{(\text{arr})}(n) \le \tau\,\bigr\},
\end{equation}
the total number of audio samples deposited into the jitter buffer by
receiver-clock time~$\tau$.  This is a right-continuous, non-decreasing,
pure-jump process with jumps of size~$N$ at each packet arrival.

%% ──────────────────────────────────────────────────────────────────────────────
\subsection{Warp Rate}\label{sec:warp-rate}

\begin{definition}[Warp rate]\label{def:warp}
  The \emph{warp rate} $u(\tau) \in \R$ is the fractional adjustment
  to the nominal playback rate.  The receiver's \emph{desired
  consumption rate} of audio samples is
  \begin{equation}\label{eq:playback-rate}
    r(\tau) \;=\; F_s\bigl(1 + u(\tau)\bigr)
  \end{equation}
  samples per unit of receiver-clock time.
  When $u = 0$, playback proceeds at the nominal sample rate~$F_s$.
  In the implementation, the quantity $1 + u(\tau)$ corresponds to
  \texttt{freq\_coeff}.
\end{definition}

%% ──────────────────────────────────────────────────────────────────────────────
\subsection{Queue Dynamics}\label{sec:queue}

The jitter buffer holds samples that have arrived but have not yet
been played back.  Samples are added in bursts of~$N$ at each packet
arrival (via~$A$) and consumed continuously at the desired rate~$r(\tau)$,
subject to the constraint that the buffer cannot go negative: when the
buffer is empty, silence is inserted instead of audio.

\begin{definition}[Queue depth and underrun regulator]\label{def:queue}
  The \emph{queue depth} $q(\tau)$ and the \emph{underrun regulator}
  $L(\tau)$ are the unique pair of right-continuous processes satisfying
  the following system for all $\tau \ge 0$:
  \begin{equation}\label{eq:queue}
    q(\tau) \;=\; q(0) \;+\; A(\tau)
             \;-\; \int_0^{\tau} F_s\bigl(1 + u(s)\bigr)\,\dd s
             \;+\; L(\tau),
  \end{equation}
  with the conditions
  \begin{enumerate}[label=(\roman*)]
    \item \textbf{Non-negativity.}\; $q(\tau) \ge 0$ for all
      $\tau \ge 0$.
    \item \textbf{Monotonicity.}\; $L$ is non-decreasing and
      $L(0) = 0$.
    \item \textbf{Complementarity.}\; $L$ increases only when $q = 0$:
      \begin{equation}\label{eq:complementarity}
        \int_0^{\infty} \mathbf{1}\{q(\tau) > 0\}\;\dd L(\tau) \;=\; 0.
      \end{equation}
  \end{enumerate}
\end{definition}

Equation~\eqref{eq:queue} is a one-sided Skorokhod reflection of the
\emph{free queue process}
\begin{equation}\label{eq:free-queue}
  \tilde{q}(\tau) \;=\; q(0) + A(\tau)
    - \int_0^{\tau} F_s\bigl(1 + u(s)\bigr)\,\dd s,
\end{equation}
which is the queue depth that would result if playback were never
interrupted by an empty buffer.
The regulator $L$ is the minimal non-negative correction that keeps
$q \ge 0$; it is given explicitly by
\begin{equation}\label{eq:regulator-explicit}
  L(\tau) \;=\; \sup_{0 \le s \le \tau}
    \bigl(-\tilde{q}(s)\bigr)^{\!+},
\end{equation}
where $x^+ = \max(x, 0)$.

\begin{definition}[Underrun]\label{def:underrun}
  An \emph{underrun event} occurs at time~$\tau$ when $q(\tau) = 0$
  and the desired consumption rate $r(\tau) > 0$.
  During an underrun, the receiver outputs silence.

  The regulator $L(\tau)$ has a direct physical interpretation:
  it is the cumulative number of \emph{silence samples} inserted
  up to time~$\tau$ in place of audio that was not available.
  The \emph{underrun time} is
  \begin{equation}\label{eq:underrun-time}
    U(\tau) \;=\; \int_0^{\tau} \mathbf{1}\{q(s) = 0\}\,\dd s,
  \end{equation}
  the total duration for which the buffer has been empty.
\end{definition}

\begin{remark}[Relationship between $L$ and $U$]
  The processes $L$ and $U$ are related but distinct.  The underrun
  time~$U$ measures duration; the regulator~$L$ measures samples.
  By the complementarity condition, $L$ increases only on the set
  $\{q = 0\}$, so
  \begin{equation}\label{eq:L-U-relation}
    L(\tau) \;=\; \int_0^{\tau} F_s\bigl(1 + u(s)\bigr)
      \,\mathbf{1}\{q(s) = 0\}\,\dd s.
  \end{equation}
\end{remark}

\begin{definition}[Overrun and session restart]\label{def:overrun}
  The system enforces an upper bound $Q_{\max}$ on the queue depth.
  The \emph{overrun stopping time} is
  \begin{equation}\label{eq:overrun-time}
    T_{\text{over}} \;=\; \inf\bigl\{\,\tau \ge 0 \;:\;
      q(\tau) > Q_{\max}\,\bigr\},
  \end{equation}
  with $T_{\text{over}} = \infty$ if the bound is never exceeded.
  When $T_{\text{over}} < \infty$, the receiver terminates the
  current session and initiates a \emph{session restart}: the buffer
  is flushed, all estimator state is reset, and a new session begins
  with fresh initial conditions.

  The queue dynamics \eqref{eq:queue}--\eqref{eq:complementarity}
  and the controller \eqref{eq:controller} are therefore only defined
  on $[0, T_{\text{over}})$.  A well-designed controller should ensure
  $T_{\text{over}} = \infty$ almost surely under normal operating
  conditions.
\end{definition}

%% ──────────────────────────────────────────────────────────────────────────────
\subsection{Observable Information}\label{sec:observables}

The receiver does not have access to the absolute reference time~$t$,
the sender clock~$\tau_A$, or the drift parameters $\alpha$, $\beta$,
$\delta$ individually.

At receiver-clock time~$\tau$, the following quantities are observable:
\begin{enumerate}[label=(\alph*)]
  \item The queue depth $q(s)$ for $s \le \tau$.
  \item The arrival times $\tau_B^{(\text{arr})}(n)$ for all packets
    received by time~$\tau$.
  \item The RTP timestamps carried by each received packet, which
    encode the sender-clock emission times
    $\tau_A^{(\text{send})}(n) = n\,\Delta$.
\end{enumerate}

We collect these into the \emph{observable filtration}
$\{\mathcal{F}_\tau\}_{\tau \ge 0}$, the $\sigma$-algebra generated
by the history of~(a)--(c) up to time~$\tau$.

%% ──────────────────────────────────────────────────────────────────────────────
\subsection{Definition of a Controller}\label{sec:controller-def}

\begin{definition}[Controller]\label{def:controller}
  A \emph{controller} is a causal (non-anticipative) functional
  \begin{equation}\label{eq:controller}
    u(\tau) \;=\; \Gamma\!\bigl(\mathcal{F}_\tau\bigr),
  \end{equation}
  mapping the observable history up to time~$\tau$ to a warp
  rate~$u(\tau) \in \R$.
  The closed-loop system is the coupled solution of
  \eqref{eq:queue}--\eqref{eq:complementarity} and
  \eqref{eq:controller}.
\end{definition}

%% ══════════════════════════════════════════════════════════════════════════════
\section{Objectives}\label{sec:objectives}

A controller $\Gamma$ induces a closed-loop trajectory
$(q, u, L, U)$ on $[0, T_{\text{over}})$.
We now formalise the criteria by which controllers are compared.

%% ──────────────────────────────────────────────────────────────────────────────
\subsection{Warp Smoothness}\label{sec:warp-smooth}

Abrupt changes in the warp rate $u$ produce audible pitch
transients.  A constant warp rate (matching the true drift)
is inaudible; only its \emph{variation} causes artefacts.
We therefore penalise the temporal roughness of $u$.

\begin{definition}[Warp cost]\label{def:warp-cost}
  For a controller producing warp trajectory $u$ on $[0, T]$,
  the \emph{warp cost} is the $\dot{H}^1$ seminorm
  \begin{equation}\label{eq:warp-cost}
    J_u(T) \;=\; \int_0^{T}\!
      \left|\frac{\dd u}{\dd \tau}(\tau)\right|^2 \dd\tau.
  \end{equation}
  A trajectory with $J_u = 0$ corresponds to a constant warp rate.
\end{definition}

\begin{remark}
  The $\dot{H}^1$ seminorm penalises only the derivative of~$u$,
  not its magnitude.  This is deliberate: the steady-state value
  of~$u$ is dictated by the clock drift~$\delta$ and is not a
  design choice.  What the controller \emph{can} control is how
  smoothly it converges to that value.
\end{remark}

%% ──────────────────────────────────────────────────────────────────────────────
\subsection{Latency Regulation}\label{sec:latency-obj}

Let $q^*$ denote the \emph{target queue depth} (in samples) and
define the \emph{queue error}
\begin{equation}\label{eq:queue-error}
  e(\tau) \;=\; q(\tau) - q^*.
\end{equation}

We decompose the latency performance into two components.

\paragraph{Bias (mean error).}
The long-run time-averaged queue error is
\begin{equation}\label{eq:bias}
  B(T) \;=\; \frac{1}{T}\int_0^{T} e(\tau)\,\dd\tau.
\end{equation}
A well-tuned controller should drive $B(T) \to 0$ as $T \to \infty$.

\paragraph{Variance (error fluctuation).}
The time-averaged mean-square error is
\begin{equation}\label{eq:mse}
  V(T) \;=\; \frac{1}{T}\int_0^{T} e(\tau)^2\,\dd\tau.
\end{equation}
This captures both persistent bias and stochastic fluctuation.
In the ergodic regime where $B(T) \to 0$,
$V(T)$ converges to the stationary variance of the queue error.

\begin{remark}[Multi-speaker synchronisation]\label{rem:multi-speaker}
  In multi-speaker setups, several receivers play the same audio
  stream.  Each receiver~$i$ runs its own controller and produces
  its own queue trajectory $q_i(\tau)$ with error $e_i(\tau)$.
  Phase coherence across speakers requires that the
  \emph{inter-speaker} variance
  \begin{equation}\label{eq:inter-speaker}
    V_{\text{sync}}(\tau)
      \;=\; \frac{1}{K}\sum_{i=1}^{K}
        \bigl(q_i(\tau) - \bar{q}(\tau)\bigr)^2,
    \qquad
    \bar{q}(\tau) = \frac{1}{K}\sum_{i=1}^{K} q_i(\tau),
  \end{equation}
  be small.  Since each $q_i$ fluctuates independently around its
  target, minimising the per-receiver variance~$V$ directly
  minimises $V_{\text{sync}}$.
\end{remark}

%% ──────────────────────────────────────────────────────────────────────────────
\subsection{Underrun and Overrun Avoidance}\label{sec:avoid-obj}

\subsubsection{Overruns}

An overrun triggers a session restart (Def.~\ref{def:overrun}),
which is a catastrophic event: all accumulated state is lost and
the stream audibly glitches.  Since the process terminates at
$T_{\text{over}}$, overrun avoidance is naturally expressed via the
stopping time.

\begin{enumerate}[label=(O\arabic*)]
  \item \textbf{Survival probability.}\;
    For a given time horizon~$T$,
    \begin{equation}\label{eq:overrun-survival}
      P\bigl(T_{\text{over}} > T\bigr) \;\to\; 1.
    \end{equation}
  \item \textbf{Expected lifetime.}\;
    $\E[T_{\text{over}}] \to \infty$, or more precisely, requiring
    that all moments $\E[T_{\text{over}}^k]$ grow appropriately.
\end{enumerate}

A well-designed controller should ensure
$T_{\text{over}} = \infty$ a.s.\ under the stationary drift
and jitter assumptions
(\cref{as:drift,as:jitter}).

\subsubsection{Underruns}

Unlike overruns, an underrun does not terminate the process---the
receiver inserts silence and continues.  A good controller should
recover from underruns and prevent them from recurring.  This
suggests metrics that measure the \emph{density} of underruns rather
than a single stopping time.

\begin{enumerate}[label=(U\arabic*)]
  \item \textbf{Underrun rate (samples).}\;
    The long-run average rate of silence insertion:
    \begin{equation}\label{eq:underrun-rate}
      R_L \;=\; \limsup_{T \to \infty}\;
        \frac{L(T)}{T}.
    \end{equation}
    This has units of samples per unit time and should be driven
    to zero.

  \item \textbf{Underrun fraction (time).}\;
    The fraction of time spent in underrun:
    \begin{equation}\label{eq:underrun-fraction}
      R_U \;=\; \limsup_{T \to \infty}\;
        \frac{U(T)}{T}.
    \end{equation}
    This is the long-run probability that the buffer is empty at a
    uniformly chosen instant.

  \item \textbf{Underrun event rate.}\;
    Let $N_{\text{under}}(T)$ count the number of transitions
    $q(\tau) : {>}0 \;\to\; 0$ on $[0, T]$.  The event rate
    \begin{equation}\label{eq:underrun-event-rate}
      R_N \;=\; \limsup_{T \to \infty}\;
        \frac{N_{\text{under}}(T)}{T}
    \end{equation}
    measures how frequently underruns occur, regardless of their
    duration.
\end{enumerate}

\begin{remark}
  These three metrics are not equivalent.  A single long underrun
  contributes heavily to $R_L$ and $R_U$ but only once to $R_N$.
  Conversely, many brief underruns inflate $R_N$ while contributing
  little to $R_U$.  In practice, all three should be small; the
  choice of which to optimise depends on the application's
  sensitivity to silence gaps versus playback interruptions.
\end{remark}

%% ══════════════════════════════════════════════════════════════════════════════
\section{Fluid Approximation}\label{sec:fluid}

The exact queue dynamics (\cref{def:queue}) involve a pure-jump
arrival process, making direct transfer-function analysis difficult.
In this section we replace the discrete packet arrivals with a
continuous flow and derive an It\^{o} stochastic differential equation
for the queue error.  All quantities remain in receiver-clock
time~$\tau$.

%% ──────────────────────────────────────────────────────────────────────────────
\subsection{Discrete Queue Recurrence}\label{sec:discrete-recurrence}

We first derive the exact per-packet queue dynamics.  Let
\begin{equation}\label{eq:q-sampled}
  q_n \;\coloneqq\; q\!\bigl(\tau_B^{(\text{arr})}(n)^{+}\bigr)
\end{equation}
denote the queue depth immediately after the $n$-th packet's $N$~samples
have been deposited.  Between arrivals $n{-}1$ and $n$, the receiver
consumes samples at rate $F_s(1+u)$ for the inter-arrival
duration
\begin{equation}\label{eq:iat-recall}
  \Delta\tau(n)
    \;\coloneqq\; \tau_B^{(\text{arr})}(n) - \tau_B^{(\text{arr})}(n-1)
    \;=\; \frac{\Delta}{1+\delta} + (1+\beta)\,\xi(n)
\end{equation}
(from \eqref{eq:inter-arrival}).
Ignoring the reflecting barrier (i.e.\ assuming $q > 0$ throughout),
the recurrence is
\begin{equation}\label{eq:q-recurrence}
  q_n \;=\; q_{n-1} \;+\; N
        \;-\; F_s\bigl(1 + u_n\bigr)\,\Delta\tau(n),
\end{equation}
where $u_n$ is the warp rate in effect during the interval
(assumed constant over one packet period).

\paragraph{Expanding.}
Substituting $N = F_s\,\Delta$ and
\eqref{eq:iat-recall} into \eqref{eq:q-recurrence}:
\begin{align}
  q_n - q_{n-1}
    &= F_s\,\Delta
       \;-\; F_s(1+u_n)\!\left[
         \frac{\Delta}{1+\delta} + (1+\beta)\,\xi(n)\right]
       \notag \\[4pt]
    &= F_s\,\Delta\!\left[1 - \frac{1+u_n}{1+\delta}\right]
       \;-\; F_s(1+u_n)(1+\beta)\,\xi(n).
       \label{eq:dq-exact}
\end{align}

The first term is the deterministic drift. Since
$1 - \frac{1+u}{1+\delta} = \frac{\delta - u}{1+\delta}$,
it equals
\begin{equation}\label{eq:dq-drift}
  F_s\,\Delta\;\frac{\delta - u_n}{1+\delta}.
\end{equation}

The second term is the stochastic jitter contribution:
\begin{equation}\label{eq:dq-jitter}
  -\,F_s(1+u_n)(1+\beta)\,\xi(n).
\end{equation}

%% ──────────────────────────────────────────────────────────────────────────────
\subsection{Summing Over a Control Epoch}\label{sec:epoch-sum}

Consider a control epoch of receiver-clock duration~$h$, during which
the warp rate $u$ is held constant.  The number of packets arriving
in this epoch is
\begin{equation}\label{eq:epoch-count}
  M(h) \;=\; \#\bigl\{\,n : \tau_B^{(\text{arr})}(n) \in
    [\tau,\,\tau + h]\,\bigr\}.
\end{equation}
Since $\E[\Delta\tau(n)] = \Delta/(1+\delta)$, the expected count is
\begin{equation}\label{eq:epoch-count-mean}
  \E[M(h)] \;=\; \frac{h(1+\delta)}{\Delta}.
\end{equation}

Summing \eqref{eq:dq-exact} over the $M(h)$ packets in the epoch:
\begin{equation}\label{eq:dq-epoch}
  q(\tau+h) - q(\tau)
    \;=\; F_s(\delta - u)\,h
    \;-\; F_s(1+u)(1+\beta)\sum_{n \in \text{epoch}} \xi(n),
\end{equation}
where we have used
$M(h) \cdot \frac{F_s\,\Delta}{1+\delta} = F_s\,h$
for the deterministic term (replacing $M$ by its mean; the
fluctuation in $M$ is itself a jitter effect already captured
by the $\xi$~sum).

%% ──────────────────────────────────────────────────────────────────────────────
\subsection{Central Limit Passage}\label{sec:clt}

\begin{assumption}[Timescale separation]\label{as:timescale}
  The control update period~$h$ satisfies $h \gg \Delta$ and
  $h \gg M\,\Delta$, where $M$ is the mixing horizon from
  \cref{as:jitter}.  That is, many packets arrive per control
  epoch, and the epoch spans many mixing lengths.
\end{assumption}

Under \cref{as:timescale}, the epoch contains $M(h) \gg 1$ jitter
increments $\xi(n)$, many of which are approximately independent
(by the mixing condition in \cref{as:jitter}).  By Donsker's
invariance principle for mixing sequences, the partial-sum process
\begin{equation}\label{eq:partial-sum}
  S(m) \;=\; \sum_{n=1}^{m} \xi(n)
\end{equation}
converges (in distribution, under diffusive rescaling) to a Brownian
motion:
\begin{equation}\label{eq:donsker}
  \frac{1}{\sqrt{m}}\,S(\lfloor m\,\cdot \rfloor)
  \;\xRightarrow{\;d\;}\; \sigma_{\text{eff}}\;W(\cdot),
  \qquad m \to \infty,
\end{equation}
where $W$ is a standard Wiener process and the
\emph{effective jitter standard deviation} is
\begin{equation}\label{eq:sigma-eff}
  \sigma_{\text{eff}}^2
    \;=\; \sum_{k=-\infty}^{\infty}
      \mathrm{Cov}\bigl(\xi(0),\,\xi(k)\bigr)
    \;=\; \sigma_\xi^2 + 2\sum_{k=1}^{\infty}
      \mathrm{Cov}\bigl(\xi(0),\,\xi(k)\bigr).
\end{equation}
This is the spectral density of the $\xi$~process evaluated at
frequency zero.  If $\xi$ is uncorrelated (i.i.d.), then
$\sigma_{\text{eff}} = \sigma_\xi$.  Short-range positive
correlations increase $\sigma_{\text{eff}}$; negative correlations
decrease it.

\begin{remark}
  The convergence
  $\sigma_{\text{eff}}^2 < \infty$ is guaranteed by the mixing
  condition in \cref{as:jitter}: the autocovariances decay fast
  enough for the series to converge absolutely.
\end{remark}

%% ──────────────────────────────────────────────────────────────────────────────
\subsection{The Fluid SDE}\label{sec:fluid-sde}

Applying the CLT result to \eqref{eq:dq-epoch}, the queue
depth~$q(\tau)$ satisfies the It\^{o} stochastic differential equation
\begin{equation}\label{eq:fluid-sde}
  \boxed{%
    \dd q \;=\; F_s(\delta - u)\,\dd\tau
          \;+\; \sigma_w\,\dd W(\tau),
  }
\end{equation}
where $W$ is a standard Wiener process and the
\emph{diffusion coefficient} is
\begin{equation}\label{eq:sigma-w}
  \sigma_w
    \;=\; F_s\,(1+u)(1+\beta)\;\sigma_{\text{eff}}\,
          \sqrt{\frac{1+\delta}{\Delta}}.
\end{equation}

\paragraph{Derivation of $\sigma_w$.}
In one epoch of duration $h$, the noise term in \eqref{eq:dq-epoch}
has variance
\begin{equation}\label{eq:noise-var}
  \Var\!\left[F_s(1+u)(1+\beta)\sum_{n=1}^{M(h)} \xi(n)\right]
  \;=\; F_s^2(1+u)^2(1+\beta)^2 \;\cdot\; M(h)\,\sigma_{\text{eff}}^2.
\end{equation}
Substituting $M(h) = h(1+\delta)/\Delta$ from
\eqref{eq:epoch-count-mean}, the variance is proportional to~$h$:
\begin{equation}\label{eq:noise-var-h}
  \sigma_w^2\,h
  \;=\; F_s^2(1+u)^2(1+\beta)^2\,
        \frac{(1+\delta)}{\Delta}\,
        \sigma_{\text{eff}}^2\;h,
\end{equation}
confirming the diffusive scaling and yielding \eqref{eq:sigma-w}.

\paragraph{Simplified form.}
Since $|u|, |\beta|, |\delta| \ll 1$ (all at the ppm level), to
leading order
\begin{equation}\label{eq:sigma-w-simple}
  \sigma_w \;=\; \frac{F_s\,\sigma_{\text{eff}}}{\sqrt{\Delta}}
           \;=\; \frac{F_s^{3/2}\,\sigma_{\text{eff}}}{\sqrt{N}}.
\end{equation}
We use this simplified form in subsequent analysis.

\begin{remark}
  The second equality follows from $\Delta = N/F_s$.
  The dependence $\sigma_w \propto 1/\sqrt{N}$ is intuitive: larger
  packets (more samples per packet) mean fewer arrivals per unit time,
  so each jitter increment has more influence per unit of consumed time.
\end{remark}

%% ──────────────────────────────────────────────────────────────────────────────
\subsection{Queue Error SDE}\label{sec:error-sde}

In terms of the queue error $e(\tau) = q(\tau) - q^*$
(\cref{eq:queue-error}), the fluid SDE becomes
\begin{equation}\label{eq:error-sde}
  \dd e \;=\; F_s(\delta - u)\,\dd\tau
        \;+\; \sigma_w\,\dd W(\tau),
\end{equation}
since $q^*$ is constant.  This is valid on the region $\{q > 0\}
\cap \{q < Q_{\max}\}$, i.e.\ away from the reflecting barrier and
the overrun boundary.

%% ──────────────────────────────────────────────────────────────────────────────
\subsection{Validity and Limitations}\label{sec:fluid-validity}

The fluid approximation \eqref{eq:fluid-sde} is accurate when:
\begin{enumerate}[label=(\alph*)]
  \item The control timescale is much longer than the packet
    period~$\Delta$ (\cref{as:timescale}).
  \item The queue depth~$q$ is well above zero, so the reflecting
    barrier is not active.  Near $q = 0$, the discrete nature of
    arrivals matters: an underrun lasts until the next packet
    arrives (a random time on the order of~$\Delta$), which the
    continuous model cannot capture.
  \item The jitter increments $\xi(n)$ have finite
    $\sigma_{\text{eff}}$ (\cref{eq:sigma-eff}).  Heavy-tailed
    jitter (infinite variance) would invalidate the CLT and
    require a different scaling limit (e.g.\ $\alpha$-stable
    processes).
\end{enumerate}

%% ══════════════════════════════════════════════════════════════════════════════
\section{Analysis of the Proportional Controller}\label{sec:prop-analysis}

We now analyse the simplest non-trivial controller under the fluid
model: feedforward drift cancellation plus proportional feedback on the
queue error.  All results in this section follow from elementary
properties of the Ornstein--Uhlenbeck process.

%% ──────────────────────────────────────────────────────────────────────────────
\subsection{Pure Feedforward is Insufficient}\label{sec:ff-fail}

Suppose the controller has perfect knowledge of the relative
drift~$\delta$ and sets $u(\tau) = \delta$ for all~$\tau$ (pure
feedforward, no feedback).  The error SDE
\eqref{eq:error-sde} becomes
\begin{equation}\label{eq:ff-sde}
  \dd e \;=\; \sigma_w\,\dd W(\tau),
\end{equation}
a driftless Brownian motion.  Its variance grows without bound:
\begin{equation}\label{eq:ff-var}
  \Var[e(\tau)] \;=\; \sigma_w^2\;\tau.
\end{equation}
By the recurrence of one-dimensional Brownian motion, the process
hits every level in finite time almost surely.  In particular,
both the underrun barrier ($e = -q^*$) and the overrun barrier
($e = Q_{\max} - q^*$) are reached in finite expected time:
\begin{equation}\label{eq:ff-hitting}
  \E[T_{\text{hit}}] \;=\; \frac{q^*\,(Q_{\max} - q^*)}{\sigma_w^2}.
\end{equation}
No amount of buffering can prevent eventual session death under
pure feedforward control.

%% ──────────────────────────────────────────────────────────────────────────────
\subsection{Feedforward Plus Proportional Feedback}\label{sec:ff-prop}

Now consider the controller
\begin{equation}\label{eq:prop-controller}
  u(\tau) \;=\; \delta + K\,e(\tau),
\end{equation}
where $K > 0$ is the proportional gain.  Substituting into
\eqref{eq:error-sde}:
\begin{equation}\label{eq:ou-sde}
  \dd e \;=\; -F_s\,K\,e\;\dd\tau \;+\; \sigma_w\,\dd W(\tau).
\end{equation}
This is the \emph{Ornstein--Uhlenbeck} (OU) SDE with
mean-reversion rate $\theta = F_s\,K$ and diffusion
coefficient~$\sigma_w$.

%% ──────────────────────────────────────────────────────────────────────────────
\subsection{Stationary Distribution}\label{sec:ou-stationary}

The OU process \eqref{eq:ou-sde} has a unique stationary
distribution:
\begin{equation}\label{eq:ou-stat}
  e \;\sim\; \mathcal{N}\!\left(0,\;\frac{\sigma_w^2}{2\,F_s\,K}\right)
  \qquad\text{(in steady state)}.
\end{equation}
Unlike the pure feedforward case, the variance is bounded and
time-independent.

%% ──────────────────────────────────────────────────────────────────────────────
\subsection{The Fundamental Tradeoff}\label{sec:tradeoff}

The controller gain $K$ simultaneously determines the distributions
of the queue error~$e$ and the warp rate~$u$.  Since
$u = \delta + K\,e$ and $\delta$ is (approximately) constant,
the stochastic part of~$u$ is $K\,e$.

\begin{proposition}[Stationary variances]\label{prop:variances}
  Under the proportional controller \eqref{eq:prop-controller}
  with the fluid model \eqref{eq:error-sde}, the stationary
  variances of the queue error and the warp rate are
  \begin{align}
    \Var[e] &\;=\; \frac{\sigma_w^2}{2\,F_s\,K},
      \label{eq:var-e} \\[6pt]
    \Var[u] &\;=\; K^2\,\Var[e]
             \;=\; \frac{K\,\sigma_w^2}{2\,F_s}.
      \label{eq:var-u}
  \end{align}
\end{proposition}

These move in \textbf{opposite directions} as $K$ varies:
\begin{center}
\begin{tabular}{@{}lcc@{}}
  \toprule
  & $K \uparrow$ & $K \downarrow$ \\
  \midrule
  $\Var[e]$ (queue error) & $\downarrow$ & $\uparrow$ \\
  $\Var[u]$ (warp variation) & $\uparrow$ & $\downarrow$ \\
  \bottomrule
\end{tabular}
\end{center}

\noindent
Eliminating~$K$ between \eqref{eq:var-e} and \eqref{eq:var-u} yields
the \emph{Pareto frontier}:
\begin{equation}\label{eq:pareto}
  \boxed{\;\Var[e] \;\cdot\; \Var[u]
    \;=\; \frac{\sigma_w^4}{4\,F_s^2}.\;}
\end{equation}
This is a hyperbola in the $(\Var[e],\,\Var[u])$ plane.  No
proportional controller can achieve a point below this curve: any
reduction in queue error variance must be paid for by an increase in
warp rate variance, and vice versa.  The product is fixed by the
jitter intensity~$\sigma_w$ and the sample rate~$F_s$.

\begin{remark}
  The Pareto frontier \eqref{eq:pareto} arises from the
  \emph{proportional} controller family.
  A more general controller (e.g.\ one that
  also acts on $\dd u/\dd\tau$) could potentially achieve points
  below this curve, at the cost of introducing higher-order dynamics.
\end{remark}

%% ──────────────────────────────────────────────────────────────────────────────
\subsection{Underrun Probability}\label{sec:underrun-prob}

In steady state, the queue error $e$ is Gaussian with zero mean and
variance $\Var[e]$ from \eqref{eq:var-e}.  An underrun occurs when
$e < -q^*$, i.e.\ $q < 0$.  The steady-state probability of being
in an underrun state is
\begin{equation}\label{eq:underrun-tail}
  P(q < 0)
    \;=\; P(e < -q^*)
    \;=\; \Phi\!\left(\frac{-q^*}{\sqrt{\Var[e]}}\right)
    \;=\; \Phi\!\left(-q^*\,\sqrt{\frac{2\,F_s\,K}{\sigma_w^2}}\right),
\end{equation}
where $\Phi$ is the standard normal CDF.  This decreases rapidly
(Gaussian tail) as either $q^*$ or $K$ increases.

\begin{remark}
  This is the steady-state occupancy probability
  $R_U$ from \cref{eq:underrun-fraction}.  For practical targets
  (e.g.\ $q^* = 5\sqrt{\Var[e]}$), the underrun probability is
  negligibly small ($< 10^{-6}$).  The tradeoff is that a larger
  $q^*$ means higher nominal latency.
\end{remark}

%% ══════════════════════════════════════════════════════════════════════════════
\section{Spectral Analysis}\label{sec:spectral}

The results of \cref{sec:prop-analysis} were derived via the
time-domain properties of the OU process.  In this section we
re-derive them in the frequency domain, which reveals a more general
principle: for linear systems driven by Gaussian noise, the
\emph{transfer function completely determines the output probability
distribution}.

%% ──────────────────────────────────────────────────────────────────────────────
\subsection{Plant Transfer Function}\label{sec:plant-tf}

The fluid error SDE \eqref{eq:error-sde}, written in input--output
form, is
\begin{equation}\label{eq:plant-io}
  \dot{e}(\tau) \;=\; -F_s\,u(\tau) \;+\; w(\tau),
\end{equation}
where $w(\tau) = \sigma_w\,\dot{W}(\tau)$ is white Gaussian noise
with (two-sided) power spectral density
\begin{equation}\label{eq:psd-input}
  S_w(\omega) \;=\; \sigma_w^2
\end{equation}
and we have absorbed the constant $F_s\,\delta$ into the feedforward
part of~$u$ (i.e.\ we take $u$ to include the drift cancellation,
so the remaining dynamics are driven by noise alone).

Taking the Laplace transform of \eqref{eq:plant-io} with zero
initial conditions:
\begin{equation}\label{eq:plant-laplace}
  s\,E(s) \;=\; -F_s\,U(s) \;+\; W(s).
\end{equation}
The open-loop plant from $U$ to $E$ is the pure integrator
\begin{equation}\label{eq:plant-G}
  G(s) \;=\; \frac{E(s)}{W(s) - F_s\,U(s)}
       \;=\; \frac{1}{s},
\end{equation}
and the noise enters the error through the same integrator:
$E(s) = \frac{1}{s}\bigl[W(s) - F_s\,U(s)\bigr]$.

\begin{remark}[Intuition: the integrator]
  The plant $G(s) = 1/s$ is an integrator: it accumulates its input
  over time.  In the frequency domain, this means that slow
  disturbances (low~$\omega$) are amplified much more than fast ones
  (high~$\omega$), since $|1/j\omega|$ is large when $\omega$ is small.
  Physically, a slow persistent jitter bias drains or fills the buffer
  steadily, while a fast oscillating jitter averages out before the
  queue can respond.  This is why drift (a zero-frequency disturbance)
  is so dangerous: the integrator amplifies it to infinity.
\end{remark}

%% ──────────────────────────────────────────────────────────────────────────────
\subsection{PSD Propagation}\label{sec:psd-prop}

For any linear time-invariant system with transfer function $H(s)$
driven by a wide-sense stationary input with PSD $S_{\text{in}}(\omega)$,
the output PSD is
\begin{equation}\label{eq:psd-general}
  S_{\text{out}}(\omega) \;=\; \bigl|H(j\omega)\bigr|^2\;
    S_{\text{in}}(\omega),
\end{equation}
and the output variance (the total power) is
\begin{equation}\label{eq:var-psd}
  \Var[\text{output}]
    \;=\; \int_{-\infty}^{\infty}
      S_{\text{out}}(\omega)\;\frac{\dd\omega}{2\pi}
    \;=\; S_{\text{in}} \cdot \|H\|_{\mathcal{H}_2}^2,
\end{equation}
where $\|H\|_{\mathcal{H}_2}^2 = \int |H(j\omega)|^2\,
\frac{\dd\omega}{2\pi}$ is the squared $\mathcal{H}_2$ norm of
the transfer function.

\begin{remark}[Intuition: PSD as a frequency-by-frequency budget]
  Think of any signal as a superposition of sine waves at different
  frequencies.  The PSD $S(\omega)$ tells you how much power is
  carried by each frequency~$\omega$---it is the signal's
  ``energy budget'' across frequencies.

  A linear system cannot create new frequencies; it can only amplify
  or attenuate existing ones.  The transfer function $H(j\omega)$
  gives the complex gain at each frequency, so the output power at
  frequency~$\omega$ is the input power times $|H(j\omega)|^2$.
  The total output variance is the sum (integral) over all
  frequencies.

  The $\mathcal{H}_2$ norm $\|H\|_{\mathcal{H}_2}$ is therefore
  a single number that captures \emph{how much total noise the
  system lets through}.  A controller that minimises the
  $\mathcal{H}_2$ norm of the closed-loop transfer function is
  one that suppresses as much noise as possible.
\end{remark}

\begin{remark}[Gaussianity]
  When the input is Gaussian and the system is linear, the output
  is also Gaussian.  A Gaussian distribution is uniquely determined
  by its mean and variance.  Therefore the transfer function
  $H(s)$, together with the input PSD, \emph{completely determines
  the output probability distribution}---not merely its second
  moment.
\end{remark}

%% ──────────────────────────────────────────────────────────────────────────────
\subsection{Closed-Loop Transfer Functions}\label{sec:cl-tf}

Under the proportional controller $U(s) = K\,E(s)$, the
closed-loop becomes
\begin{equation}\label{eq:cl-E}
  s\,E(s) = -F_s\,K\,E(s) + W(s)
  \quad\Longrightarrow\quad
  E(s) = \frac{1}{s + F_s K}\;W(s).
\end{equation}

Define the closed-loop transfer functions from the noise $W$ to
each signal of interest:

\paragraph{Noise to queue error.}
\begin{equation}\label{eq:tf-e}
  H_e(s) \;=\; \frac{E(s)}{W(s)} \;=\; \frac{1}{s + F_s K}.
\end{equation}

\paragraph{Noise to warp rate.}
Since $U = K\,E$:
\begin{equation}\label{eq:tf-u}
  H_u(s) \;=\; \frac{U(s)}{W(s)} \;=\; \frac{K}{s + F_s K}.
\end{equation}

\paragraph{Noise to warp rate derivative.}
Since $\dot{u} = K\,\dot{e}$, multiplication by $s$ in the
Laplace domain gives
\begin{equation}\label{eq:tf-v}
  H_v(s) \;=\; s\,H_u(s) \;=\; \frac{K\,s}{s + F_s K}.
\end{equation}

%% ──────────────────────────────────────────────────────────────────────────────
\subsection{Output PSDs and Variances}\label{sec:output-psd}

Since $S_w = \sigma_w^2$ (constant), the output PSDs are:

\paragraph{Queue error.}
\begin{equation}\label{eq:psd-e}
  S_e(\omega) \;=\; \frac{\sigma_w^2}{(F_s K)^2 + \omega^2}.
\end{equation}
This is a Lorentzian with half-width $F_s K$.

The half-width $\omega_c = F_s K$ is the controller's
\emph{bandwidth}: it divides the frequency axis into two regimes.
\begin{itemize}
  \item \textbf{Low frequencies} ($\omega \ll F_s K$):
    $S_e(\omega) \approx \sigma_w^2 / (F_s K)^2$.  The controller
    successfully suppresses slow queue variations---the error is
    small and nearly flat.
  \item \textbf{High frequencies} ($\omega \gg F_s K$):
    $S_e(\omega) \approx \sigma_w^2 / \omega^2$.  The controller
    is too slow to react, and the error behaves like uncontrolled
    integrated noise.
\end{itemize}
Increasing $K$ widens the bandwidth: the controller can fight
faster disturbances, reducing the total queue error variance.
But this comes at a cost, as we see next.
The variance is
\begin{equation}\label{eq:var-e-psd}
  \Var[e] \;=\; \int_{-\infty}^{\infty}
    \frac{\sigma_w^2}{(F_s K)^2 + \omega^2}\;
    \frac{\dd\omega}{2\pi}
  \;=\; \frac{\sigma_w^2}{2\,F_s\,K},
\end{equation}
recovering \eqref{eq:var-e}.

\paragraph{Warp rate.}
\begin{equation}\label{eq:psd-u}
  S_u(\omega) \;=\; \frac{K^2\,\sigma_w^2}{(F_s K)^2 + \omega^2}.
\end{equation}
The variance is
\begin{equation}\label{eq:var-u-psd}
  \Var[u] \;=\; \frac{K^2\,\sigma_w^2}{2\,F_s\,K}
          \;=\; \frac{K\,\sigma_w^2}{2\,F_s},
\end{equation}
recovering \eqref{eq:var-u}.

\paragraph{Warp rate derivative.}
\begin{equation}\label{eq:psd-v}
  S_v(\omega) \;=\; \frac{K^2\,\omega^2\,\sigma_w^2}
    {(F_s K)^2 + \omega^2}.
\end{equation}
The variance is
\begin{equation}\label{eq:var-v-psd}
  \Var[\dot{u}] \;=\; K^2\,\sigma_w^2
    \int_{-\infty}^{\infty}
    \frac{\omega^2}{(F_s K)^2 + \omega^2}\;
    \frac{\dd\omega}{2\pi}
  \;=\; +\infty.
\end{equation}
The integral diverges because the integrand approaches
$K^2\sigma_w^2$ as $|\omega| \to \infty$.

\begin{proposition}[Infinite warp roughness]\label{prop:inf-roughness}
  Under the proportional controller \eqref{eq:prop-controller}
  with the fluid model, the warp cost functional
  $J_u$ (\cref{def:warp-cost}) is infinite almost surely.
  That is, the proportional controller produces warp trajectories
  that are nowhere differentiable.
\end{proposition}

\begin{remark}[Intuition: high-frequency passthrough]
  The transfer function $H_v(s) = Ks/(s + F_s K)$ is a
  \emph{high-pass filter}.  At low frequencies it attenuates
  (the controller changes $u$ slowly in response to slow queue
  drifts), but at high frequencies $|H_v(j\omega)| \to K$---it
  passes noise straight through with gain~$K$, without any
  attenuation.

  Since white noise has equal power at all frequencies, and $H_v$
  does not roll off, the warp derivative picks up an infinite
  amount of high-frequency noise.  This is the mathematical
  expression of the fact that ``reacting proportionally to every
  instantaneous queue fluctuation produces jittery, rough control
  actions.''

  A controller that instead acts on $v = \dot{u}$ (controlling the
  rate of change of the warp) would introduce an additional $1/s$
  factor, turning the high-pass into a band-pass that decays at
  high frequency, yielding finite $\Var[\dot{u}]$.
\end{remark}

%% ══════════════════════════════════════════════════════════════════════════════
\section{Second-Order Controller}\label{sec:second-order}

The proportional controller of \cref{sec:prop-analysis} acts directly
on the warp rate~$u$, producing warp trajectories with infinite
roughness (\cref{prop:inf-roughness}).  We now promote the control
input to the \emph{derivative} of the warp rate, $v = \dot{u}$,
treating $u$ as an additional state variable.

%% ──────────────────────────────────────────────────────────────────────────────
\subsection{Augmented State Space}\label{sec:augmented-state}

The augmented system has state $\mathbf{x} = (e,\,u)^{\!\top}$
and control input $v = \dot{u}$.  Absorbing the drift feedforward
($u$ represents the feedback component after subtracting~$\delta$),
the dynamics are
\begin{equation}\label{eq:augmented-dynamics}
  \dd\!\begin{pmatrix} e \\ u \end{pmatrix}
  \;=\;
  \underbrace{%
    \begin{pmatrix} 0 & -F_s \\ 0 & 0 \end{pmatrix}
  }_{A}
  \begin{pmatrix} e \\ u \end{pmatrix}\dd\tau
  \;+\;
  \underbrace{%
    \begin{pmatrix} 0 \\ 1 \end{pmatrix}
  }_{B}
  v\;\dd\tau
  \;+\;
  \underbrace{%
    \begin{pmatrix} \sigma_w \\ 0 \end{pmatrix}
  }_{G}
  \dd W.
\end{equation}
The open-loop plant (from $v$ to $e$, without noise) is a
\emph{double integrator}: $v$ drives $u$, which in turn drives $e$
through the integrator $G(s) = 1/s$.

%% ──────────────────────────────────────────────────────────────────────────────
\subsection{Control Law}\label{sec:2nd-control-law}

\begin{definition}[Second-order controller]\label{def:2nd-controller}
  The \emph{second-order controller} sets
  \begin{equation}\label{eq:2nd-law}
    v(\tau) \;=\; K_1\,e(\tau) \;-\; K_2\,u(\tau),
  \end{equation}
  where $K_1 > 0$ and $K_2 > 0$ are gains.  The term $K_1\,e$
  provides a restoring force proportional to the queue error,
  and $-K_2\,u$ provides damping proportional to the current
  warp rate.
\end{definition}

\begin{remark}[Physical analogy: damped harmonic oscillator]
  Under \eqref{eq:2nd-law}, the queue error satisfies the
  second-order ODE (ignoring noise):
  \begin{equation}\label{eq:damped-oscillator}
    \ddot{e} + K_2\,\dot{e} + F_s\,K_1\,e \;=\; 0.
  \end{equation}
  This is a damped harmonic oscillator with natural frequency and
  damping ratio
  \begin{equation}\label{eq:oscillator-params}
    \omega_n \;=\; \sqrt{F_s\,K_1},
    \qquad
    \zeta \;=\; \frac{K_2}{2\,\omega_n}
         \;=\; \frac{K_2}{2\sqrt{F_s\,K_1}}.
  \end{equation}
  The queue error behaves like a mass on a spring, kicked by jitter
  noise, with $K_1$ controlling the spring stiffness (how fast it
  returns to centre) and $K_2$ controlling the damping (how quickly
  oscillations die out).
\end{remark}

%% ──────────────────────────────────────────────────────────────────────────────
\subsection{Closed-Loop Transfer Functions}\label{sec:2nd-tf}

In the Laplace domain, $V(s) = K_1\,E(s) - K_2\,U(s)$.
With $s\,U = V$ (i.e.\ $U = V/s$), we have
$(s + K_2)\,U = K_1\,E$, so
\begin{equation}\label{eq:2nd-U-E}
  U(s) \;=\; \frac{K_1}{s + K_2}\;E(s).
\end{equation}

Substituting into the plant equation $s\,E = -F_s\,U + W$:
\begin{equation}\label{eq:2nd-char}
  \left(s^2 + K_2\,s + F_s\,K_1\right) E(s)
    \;=\; (s + K_2)\,W(s).
\end{equation}

Denote the characteristic polynomial
\begin{equation}\label{eq:char-poly}
  D(s) \;=\; s^2 + K_2\,s + F_s\,K_1
       \;=\; s^2 + 2\zeta\omega_n\,s + \omega_n^2.
\end{equation}
All roots lie in the open left half-plane for $K_1, K_2 > 0$,
confirming stability.

The closed-loop transfer functions from the jitter noise~$W$ to
each output are:

\paragraph{Noise to queue error.}
\begin{equation}\label{eq:2nd-tf-e}
  H_e(s) \;=\; \frac{s + K_2}{D(s)}.
\end{equation}

\paragraph{Noise to warp rate.}
\begin{equation}\label{eq:2nd-tf-u}
  H_u(s) \;=\; \frac{K_1}{D(s)}.
\end{equation}

\paragraph{Noise to warp rate derivative.}
\begin{equation}\label{eq:2nd-tf-v}
  H_v(s) \;=\; s\,H_u(s) \;=\; \frac{K_1\,s}{D(s)}.
\end{equation}
This is a \emph{band-pass filter}: it rolls off as $1/s$ at high
frequency and vanishes at $s = 0$.  Unlike the first-order case
(\cref{eq:tf-v}), $|H_v(j\omega)| \to 0$ as $\omega \to \infty$.

%% ──────────────────────────────────────────────────────────────────────────────
\subsection{Stationary Variances}\label{sec:2nd-variances}

We use the standard $\mathcal{H}_2$ norm formula for second-order
systems.  For a transfer function $H(s) = (n_1\,s + n_0)/
(s^2 + a\,s + b)$ with $a, b > 0$:
\begin{equation}\label{eq:h2-formula}
  \|H\|_{\mathcal{H}_2}^2
    \;=\; \frac{n_1^2\,b \;+\; n_0^2}{2\,a\,b}.
\end{equation}

Applying this with $a = K_2$ and $b = F_s\,K_1$:

\begin{proposition}[Stationary variances, second-order controller]
  \label{prop:2nd-variances}
  Under the second-order controller \eqref{eq:2nd-law} with the
  fluid model, all three stationary variances are finite:
  \begin{align}
    \Var[e]
      &\;=\; \sigma_w^2\;\frac{F_s\,K_1 + K_2^2}{2\,K_2\,F_s\,K_1},
      \label{eq:2nd-var-e} \\[6pt]
    \Var[u]
      &\;=\; \sigma_w^2\;\frac{K_1}{2\,K_2\,F_s},
      \label{eq:2nd-var-u} \\[6pt]
    \Var[\dot{u}]
      &\;=\; \sigma_w^2\;\frac{K_1^2}{2\,K_2}.
      \label{eq:2nd-var-v}
  \end{align}
\end{proposition}

\begin{remark}[Derivation]
  For example, $H_u(s) = K_1/(s^2 + K_2\,s + F_s K_1)$ has
  $n_1 = 0$, $n_0 = K_1$.
  Applying \eqref{eq:h2-formula}:
  $\|H_u\|^2 = (0 + K_1^2)/(2 K_2 F_s K_1) = K_1/(2 K_2 F_s)$.
  Multiplying by $\sigma_w^2$ gives \eqref{eq:2nd-var-u}.
  The other cases follow analogously.
\end{remark}

\begin{remark}[Comparison with first-order controller]
  The first-order controller (\cref{sec:tradeoff}) has
  $\Var[\dot{u}] = \infty$, so the warp cost $J_u$ is undefined.
  The second-order controller makes all three quantities finite
  and provides a two-parameter family $(K_1, K_2)$ of controllers.
  The Pareto surface in the three-dimensional space
  $(\Var[e],\,\Var[u],\,\Var[\dot{u}])$ is now a smooth manifold
  parameterised by $(K_1, K_2)$, allowing principled tradeoffs
  among queue accuracy, warp magnitude, and warp smoothness.
\end{remark}

%% ──────────────────────────────────────────────────────────────────────────────
\subsection{In Terms of Oscillator Parameters}\label{sec:2nd-oscillator}

Expressing the variances in terms of the natural frequency
$\omega_n = \sqrt{F_s K_1}$ and damping ratio
$\zeta = K_2/(2\omega_n)$:
\begin{align}
  \Var[e]
    &\;=\; \frac{\sigma_w^2(1 + 4\zeta^2)}{4\,\zeta\,\omega_n},
    \label{eq:2nd-var-e-osc} \\[6pt]
  \Var[u]
    &\;=\; \frac{\sigma_w^2\,\omega_n}{4\,\zeta\,F_s^2},
    \label{eq:2nd-var-u-osc} \\[6pt]
  \Var[\dot{u}]
    &\;=\; \frac{\sigma_w^2\,\omega_n^3}{4\,\zeta\,F_s^2}.
    \label{eq:2nd-var-v-osc}
\end{align}

\begin{remark}[Role of $\omega_n$ and $\zeta$]
  The natural frequency $\omega_n$ controls the controller bandwidth:
  higher $\omega_n$ means faster response (lower $\Var[e]$, but higher
  $\Var[u]$ and $\Var[\dot{u}]$).  The damping ratio $\zeta$ controls
  oscillatory behaviour: underdamped ($\zeta < 1$) gives faster
  settling but overshoot; overdamped ($\zeta > 1$) gives slower,
  monotone convergence.  All three variances scale as $1/\zeta$,
  and critically damped ($\zeta = 1$) is not necessarily optimal.
\end{remark}

%% ──────────────────────────────────────────────────────────────────────────────
\subsection{Dropping the Warp Magnitude Objective}\label{sec:drop-var-u}

The steady-state warp rate $u$ is dominated by the
drift~$\delta$, which is on the order of $10^{-5}$--$10^{-6}$
(a few ppm).  Its stochastic variance $\Var[u]$ is even smaller
and always inaudible.  What matters for audio quality is the
\emph{rate of change} $\dot{u}$: rapid pitch changes produce
audible chirp artefacts, whereas a slowly-varying pitch offset
is imperceptible.

Moreover, controlling $\Var[\dot{u}]$ implicitly bounds
$\Var[u]$: if $\dot{u}$ is small, $u$ cannot wander far from
its steady-state value.

We therefore optimise over two objectives only:
\begin{equation}\label{eq:2obj}
  J(K_1, K_2)
    \;=\; \Var[e] \;+\; \lambda\;\Var[\dot{u}],
\end{equation}
where $\lambda > 0$ is a design parameter weighting queue accuracy
against warp smoothness.

%% ──────────────────────────────────────────────────────────────────────────────
\subsection{Optimal Gains}\label{sec:optimal-gains}

Substituting \eqref{eq:2nd-var-e} and \eqref{eq:2nd-var-v} into
\eqref{eq:2obj}:
\begin{equation}\label{eq:J-expanded}
  J \;=\; \sigma_w^2\!\left[
    \frac{1}{2\,K_2}
    + \frac{K_2}{2\,F_s\,K_1}
    + \frac{\lambda\,K_1^2}{2\,K_2}
  \right].
\end{equation}

Setting $\partial J / \partial K_1 = 0$:
\begin{equation}\label{eq:opt-K1-cond}
  -\frac{K_2}{2\,F_s\,K_1^2} + \frac{\lambda\,K_1}{K_2} \;=\; 0
  \quad\Longrightarrow\quad
  K_2^2 \;=\; 2\,\lambda\,F_s\,K_1^3.
\end{equation}

Setting $\partial J / \partial K_2 = 0$:
\begin{equation}\label{eq:opt-K2-cond}
  -\frac{1}{2\,K_2^2} + \frac{1}{2\,F_s\,K_1}
  - \frac{\lambda\,K_1^2}{2\,K_2^2} \;=\; 0
  \quad\Longrightarrow\quad
  K_2^2 \;=\; F_s\,K_1\,(1 + \lambda\,K_1^2).
\end{equation}

Equating \eqref{eq:opt-K1-cond} and \eqref{eq:opt-K2-cond}:
\begin{equation}
  2\,\lambda\,F_s\,K_1^3
    \;=\; F_s\,K_1 + \lambda\,F_s\,K_1^3
  \quad\Longrightarrow\quad
  \lambda\,K_1^2 \;=\; 1,
\end{equation}
giving the optimal gains:

\begin{proposition}[Optimal second-order gains]\label{prop:opt-gains}
  The cost $J = \Var[e] + \lambda\,\Var[\dot{u}]$ is minimised by
  \begin{equation}\label{eq:opt-K1}
    K_1^* \;=\; \frac{1}{\sqrt{\lambda}},
    \qquad
    K_2^* \;=\; \sqrt{2\,F_s\,K_1^*}
          \;=\; \sqrt{\frac{2\,F_s}{\lambda^{1/2}}}.
  \end{equation}
  The corresponding oscillator parameters are
  \begin{equation}\label{eq:opt-osc}
    \omega_n^* \;=\; \sqrt{F_s\,K_1^*}
              \;=\; \left(\frac{F_s^2}{\lambda}\right)^{\!1/4},
    \qquad
    \boxed{\;\zeta^* \;=\; \frac{1}{\sqrt{2}}
      \;\approx\; 0.707.\;}
  \end{equation}
\end{proposition}

\begin{remark}[Universal damping ratio]
  The optimal damping ratio $\zeta^* = 1/\sqrt{2}$ is
  \textbf{independent of~$\lambda$}, the jitter intensity~$\sigma_w$,
  and the sample rate~$F_s$.  Only the natural frequency~$\omega_n$
  changes with the design parameters.  The value
  $\zeta = 1/\sqrt{2}$ is the Butterworth (maximally flat magnitude)
  damping, well known in filter design and classical optimal control.
\end{remark}

%% ──────────────────────────────────────────────────────────────────────────────
\subsection{Optimal Variances}\label{sec:opt-variances}

Substituting the optimal gains into the variance expressions:
\begin{align}
  \Var[e]^*
    &\;=\; \frac{3\,\sigma_w^2}{2\sqrt{2}\;\omega_n^*},
    \label{eq:opt-var-e} \\[6pt]
  \Var[\dot{u}]^*
    &\;=\; \frac{\sigma_w^2\,(\omega_n^*)^3}{2\sqrt{2}\;F_s^2},
    \label{eq:opt-var-v} \\[6pt]
  J^*
    &\;=\; \Var[e]^* + \lambda\,\Var[\dot{u}]^*
    \;=\; \frac{3\,\sigma_w^2}{2\sqrt{2}\;\omega_n^*}
      + \frac{\sigma_w^2\,(\omega_n^*)^3}{2\sqrt{2}\;F_s^2}\cdot\lambda.
    \label{eq:opt-J}
\end{align}

Since $\omega_n^* = (F_s^2/\lambda)^{1/4}$, both terms depend
on~$\lambda$ through~$\omega_n^*$ alone: increasing $\lambda$
(more weight on smoothness) decreases $\omega_n^*$, widening
$\Var[e]$ and shrinking $\Var[\dot{u}]$.



%% ──────────────────────────────────────────────────────────────────────────────
\subsection{Steady-State Bias}\label{sec:bias}

The second-order controller has a nonzero steady-state queue error
in the presence of constant clock drift.  This section analyses the
bias, explains its root cause in terms of system type, and outlines
a path to eliminating it.

\subsubsection{Equilibrium Analysis}

In steady state all derivatives vanish.  From the plant
equation~\eqref{eq:error-sde} (deterministic part):
\begin{equation}\label{eq:ss-plant}
  \frac{\dd e}{\dd\tau} = 0
  \quad\Longrightarrow\quad
  F_s\,(\delta - u_{\text{ss}}) = 0
  \quad\Longrightarrow\quad
  u_{\text{ss}} = \delta.
\end{equation}
From the controller law~\eqref{eq:2nd-control-law}:
\begin{equation}\label{eq:ss-controller}
  \frac{\dd u}{\dd\tau} = 0
  \quad\Longrightarrow\quad
  K_1\,e_{\text{ss}} = K_2\,u_{\text{ss}} = K_2\,\delta
  \quad\Longrightarrow\quad
  \boxed{\;e_{\text{ss}} = \frac{K_2}{K_1}\,\delta.\;}
\end{equation}
The bias is proportional to the drift~$\delta$ and to the ratio
$K_2/K_1$.

\paragraph{Numerical magnitude.}
For the default gains $K_1 = 10^{-5}$ and
$K_2 = \sqrt{2\,F_s\,K_1} \approx 0.98$ at $F_s = 48\,000$:
\begin{equation}
  \frac{K_2}{K_1} \approx 9.8 \times 10^4.
\end{equation}
A drift of $\delta = 10\;\text{ppm} = 10^{-5}$ gives a steady-state
bias of
\begin{equation}
  e_{\text{ss}} \approx 0.98\;\text{samples}
  \;\approx\; 20\;\mu\text{s}\;\text{at 48\,kHz},
\end{equation}
which is negligible compared to jitter-induced fluctuations.
For larger drifts (e.g.\ $\delta = 100\;\text{ppm}$), the bias
grows to $\approx 10\;\text{samples} \approx 200\;\mu\text{s}$.

\subsubsection{Root Cause: System Type}

The open-loop transfer function from
error~$e$ to the plant input~$u$ contains exactly \emph{one}
integrator (the warp state~$u$), but the damping term
$-K_2\,u$ makes it a \emph{leaky} integrator:
\begin{equation}
  U(s) = \frac{K_1}{s + K_2}\,E(s).
\end{equation}
In classical control terminology, this is a \textbf{Type~1} system.
A Type~1 system rejects step disturbances with zero steady-state
error but has nonzero error for a ramp.  The constant drift~$\delta$
enters the plant as a step in~$\dot{e}$, which from the perspective
of the error~$e$ is a ramp input---hence the nonzero
position error.

\subsubsection{Why the H2 Cost Misses the Bias}

The H2 optimisation (Section~\ref{sec:drop-var-u}) minimises
\begin{equation}
  J = \Var[e] + \lambda\,\Var[\dot{u}],
\end{equation}
which captures only the \emph{variance} (fluctuation around the
mean), not the \emph{mean} itself.  The steady-state bias
$e_{\text{ss}} = (K_2/K_1)\,\delta$ is deterministic and does not
appear in the variance.  Consequently, the H2-optimal gains make no
attempt to reduce the bias, focusing entirely on minimising the
stochastic fluctuations.

\subsubsection{Towards Zero Bias: Type~2 Design}

To achieve $e_{\text{ss}} = 0$ for constant drift, one needs a
\textbf{Type~2} system---two integrators in the open loop.  This
requires adding a third state variable that accumulates the
integral of the error:
\begin{align}
  \frac{\dd \xi}{\dd\tau} &= e, \label{eq:type2-xi}\\[4pt]
  \frac{\dd u}{\dd\tau} &= K_1\,e + K_0\,\xi - K_2\,u,
    \label{eq:type2-u}
\end{align}
where $\xi$ is the integrated error and $K_0 > 0$ is the integral
gain.  The controller now has three gains: $K_0$ (integral),
$K_1$ (proportional), and $K_2$ (damping).

\paragraph{Steady-state verification.}
In steady state, $\dot{\xi} = 0$ implies $e_{\text{ss}} = 0$
immediately---regardless of the drift~$\delta$ or the gain
values.  The integrator state settles to
$\xi_{\text{ss}} = K_2\,\delta/K_0$, and the warp state to
$u_{\text{ss}} = \delta$ as before.

\subsubsection{PID Interpretation}

The control law~\eqref{eq:type2-u} is a PID controller applied to
the queue error:
\begin{itemize}
  \item $K_0\,\xi = K_0\!\int\!e\,\dd\tau$: \textbf{Integral}
    action (eliminates bias),
  \item $K_1\,e$: \textbf{Proportional} action (restoring force),
  \item $-K_2\,u$: \textbf{Derivative} action on the warp state
    (damping).
\end{itemize}
The previous second-order controller is the PD subset ($K_0 = 0$).

\subsubsection{Closed-Loop Transfer Functions}

The state-space representation with state $\mathbf{x} = (e, \xi, u)^\top$
and noise input $w$:
\begin{equation}
  \underbrace{
  \begin{pmatrix} \dot{e} \\ \dot{\xi} \\ \dot{u} \end{pmatrix}
  }_{\dot{\mathbf{x}}}
  =
  \underbrace{
  \begin{pmatrix}
    0 & 0 & -F_s \\
    1 & 0 & 0 \\
    K_1 & K_0 & -K_2
  \end{pmatrix}
  }_{A}
  \begin{pmatrix} e \\ \xi \\ u \end{pmatrix}
  +
  \underbrace{
  \begin{pmatrix} \sigma_w \\ 0 \\ 0 \end{pmatrix}
  }_{B}\,\dot{W}.
\end{equation}

The closed-loop characteristic polynomial is
\begin{equation}\label{eq:type2-charpoly}
  p(s) = s^3 + K_2\,s^2 + F_s\,K_1\,s + F_s\,K_0.
\end{equation}

The transfer function from noise~$w$ to queue error~$e$ is
\begin{equation}\label{eq:type2-He}
  H_e(s) = \frac{\sigma_w\,s\,(s + K_2)}{p(s)},
\end{equation}
and from noise to warp derivative $\dot{u}$:
\begin{equation}\label{eq:type2-Hu}
  H_{\dot{u}}(s) = \frac{\sigma_w\,s\,(K_1\,s + K_0)}{p(s)}.
\end{equation}

Note the \emph{extra zero at $s = 0$} in both transfer functions
compared to the PD controller.  This zero suppresses the
low-frequency response, which is precisely the mechanism by
which integral action rejects the DC disturbance (drift).

\subsubsection{Dominant-Pole Gain Design}\label{sec:dominant-pole}

Rather than solving the full 3-parameter H2 optimisation (which
yields transcendental equations), we exploit the separation of
timescales.  Choose the integral gain $K_0$ small enough that the
characteristic polynomial factors approximately as
\begin{equation}\label{eq:dominant-factor}
  p(s) \;\approx\; (s^2 + 2\zeta\omega_n\,s + \omega_n^2)\,(s + p_0),
\end{equation}
where $p_0 \ll \zeta\omega_n$ is a slow real pole governing the
integral mode.  Expanding and matching
with~\eqref{eq:type2-charpoly}:
\begin{align}
  K_2 &= 2\zeta\omega_n + p_0 \;\approx\; 2\zeta\omega_n,
    \label{eq:K2-pid}\\
  F_s\,K_1 &= \omega_n^2 + 2\zeta\omega_n\,p_0
    \;\approx\; \omega_n^2,
    \label{eq:K1-pid}\\
  F_s\,K_0 &= \omega_n^2\,p_0
    \quad\Longrightarrow\quad
    \boxed{\;K_0 = K_1\,p_0.\;}
    \label{eq:K0-pid}
\end{align}

\begin{proposition}[PID gain design]\label{prop:pid-gains}
  Keep the Butterworth-optimal PD gains
  (Proposition~\ref{prop:opt-gains}):
  \begin{equation}
    K_1 = \frac{1}{\sqrt{\lambda}},\qquad
    K_2 = \sqrt{2\,F_s\,K_1},\qquad
    \zeta = \frac{1}{\sqrt{2}}.
  \end{equation}
  Set the integral gain via its settling time $T_I$:
  \begin{equation}\label{eq:K0-from-TI}
    p_0 = \frac{1}{T_I},\qquad
    K_0 = \frac{K_1}{T_I}.
  \end{equation}
  The integral mode eliminates the steady-state bias with time
  constant~$T_I$, while the dominant oscillator dynamics are
  essentially unchanged.
\end{proposition}

\begin{remark}[Stability margin]
  The Routh condition requires $K_0 < K_1\,K_2$.  Since
  $K_0 = K_1/T_I$ and $K_2 \approx \omega_n\sqrt{2}$, the
  condition becomes $T_I > 1/K_2 \approx 1/\omega_n$.  For
  $K_1 = 10^{-5}$ and $F_s = 48\,000$: $\omega_n \approx 0.69$,
  so $T_I > 1.4\;\text{s}$.  A practical choice of
  $T_I = 30\;\text{s}$ gives a large stability margin.
\end{remark}

\subsubsection{Effect on Variances}

Solving the Lyapunov equation $AP + PA^\top + BB^\top = 0$ for the
PID state-space representation yields the exact stationary variance:
\begin{equation}\label{eq:pid-var-e}
  \Var[e]_{\text{PID}}
    = \frac{\sigma_w^2\,(K_2^2 + F_s\,K_1)}
           {2\,F_s\,(K_1\,K_2 - K_0)}.
\end{equation}
Comparing with the PD result ($K_0 = 0$):
\begin{equation}\label{eq:pid-var-ratio}
  \frac{\Var[e]_{\text{PID}}}{\Var[e]_{\text{PD}}}
    = \frac{K_1\,K_2}{K_1\,K_2 - K_0}
    = \frac{1}{1 - K_0/(K_1\,K_2)}.
\end{equation}
Since $K_0 > 0$, the PID variance is \emph{slightly larger} than
the PD variance.  The integral mode adds a slow resonance that
admits more low-frequency noise energy.

\paragraph{Numerical overhead.}
For $K_0 = K_1/T_I$ with $T_I = 30\;\text{s}$ and
$K_2 \approx 0.98$:
\begin{equation}
  \frac{K_0}{K_1\,K_2} = \frac{1}{K_2\,T_I}
    \approx \frac{1}{0.98 \times 30} \approx 3.4\%.
\end{equation}
The variance overhead is thus $\approx 3.5\%$---negligible
compared to the benefit of eliminating steady-state bias entirely.

\subsubsection{The Full Three-Objective Cost}

The complete cost function incorporating all three objectives is
\begin{equation}\label{eq:full-cost}
  J = \underbrace{\E[e]^2}_{J_1\;\text{(bias)}}
    + \underbrace{\Var[e]}_{J_2\;\text{(regulation)}}
    + \lambda\,\underbrace{\Var[\dot{u}]}_{J_3\;\text{(smoothness)}}.
\end{equation}
For the PID controller:
\begin{align}
  J_1 &= 0  &&\text{(structurally zero for any $K_0 > 0$)},\\
  J_2 &= \frac{\sigma_w^2\,(K_2^2 + F_s\,K_1)}
              {2\,F_s\,(K_1\,K_2 - K_0)}
       &&\text{(increases with $K_0$)},\\
  J_3 &\approx \frac{\sigma_w^2\,\omega_n^3}{2\sqrt{2}\,F_s^2}
       &&\text{(weakly affected)}.
\end{align}

\begin{remark}[Optimal $K_0$ depends on unknown drift]
  Since $J_1 = 0$ for any $K_0 > 0$ but $J_2$ increases with $K_0$,
  the H2-optimal integral gain is $K_0 \to 0^+$---as small as
  possible while still nonzero.  In practice, $K_0$ must be large
  enough that the integral mode settles within a reasonable time
  (e.g.\ $T_I = 30$\,s), giving the dominant-pole
  design~\eqref{eq:K0-from-TI} as the pragmatic optimum.
\end{remark}

The user controls the $J_2$--$J_3$ tradeoff through $K_1$
(equivalently $\lambda$ or $\omega_n$), while $K_0 = K_1/T_I$
handles the bias independently.

%% ──────────────────────────────────────────────────────────────────────────────
\subsection{Discrete-Time Implementation}\label{sec:implementation}

In practice the controller runs at discrete epochs of
duration~$h$ (the scaling interval).  At each epoch~$k$:

\begin{enumerate}
  \item Observe the queue error $e_k = q_k - q^*$.
  \item Update the integral state:
    \begin{equation}\label{eq:impl-xi}
      \xi_{k+1} \;=\; \xi_k \;+\; e_k\,h.
    \end{equation}
  \item Compute the control acceleration:
    \begin{equation}\label{eq:impl-v}
      v_k \;=\; K_1\,e_k \;+\; K_0\,\xi_k \;-\; K_2\,u_k.
    \end{equation}
  \item Update the warp state by Euler integration:
    \begin{equation}\label{eq:impl-u}
      u_{k+1} \;=\; u_k \;+\; v_k\,h.
    \end{equation}
  \item Clamp to safety bounds and set the resampler coefficient:
    \begin{equation}\label{eq:impl-coeff}
      \texttt{freq\_coeff}_k \;=\; 1 + \text{clamp}(u_k,\;
        -u_{\max},\; u_{\max}).
    \end{equation}
\end{enumerate}

\begin{remark}[Anti-windup]
  When the warp state~$u$ saturates at $\pm u_{\max}$, the
  integrator~$\xi$ must be frozen to prevent windup.  In the
  implementation, $\xi$ is not updated when $u$ is clamped and the
  error has the same sign as the saturation direction.  This
  ensures the integral state does not grow unboundedly during
  transients or fault conditions.
\end{remark}

\begin{remark}[No feedforward drift estimate]
  An external drift estimate $\hat{\delta}$ (e.g.\ from RTP timestamp
  ratios) is not required.  The integral action drives the
  steady-state bias to exactly zero for any constant drift, making
  feedforward unnecessary.  Omitting feedforward avoids step
  disturbances from estimator artefacts (NTP corrections, session
  restarts, bias in early samples).
\end{remark}

%% ══════════════════════════════════════════════════════════════════════════════
\section{Skewness and Asymmetric Control}%
\label{sec:skewness}

The preceding analysis assumes symmetric Gaussian noise, yielding
a symmetric stationary distribution for the queue error.  In practice,
network jitter is \emph{one-sided}: packets can be delayed but never
arrive before the minimum propagation time.  This produces a
skewed queue error distribution whose tails are asymmetric,
with direct consequences for the underrun probability.

This section develops a Fokker--Planck analysis of the queue error
under asymmetric feedback, yielding a closed-form stationary density
and exact skewness control.

%% ──────────────────────────────────────────────────────────────────────────────
\subsection{Physical Origin of Skewness}%
\label{sec:skew-origin}

Each packet's transit delay decomposes as
$d_k = d_{\min} + j_k$ where $j_k \geq 0$.  The non-negative
constraint on~$j_k$ is physical (speed of light plus serialization)
and implies that per-packet jitter has strictly positive skewness:
$\kappa_3[j_k] > 0$.  Common models include exponential jitter
($\gamma_3 = 2$) and gamma jitter ($\gamma_3 = 2/\sqrt{\alpha}$).

The aggregate noise $w_k$ driving the queue in each control epoch
inherits this skewness.  By the CLT with Edgeworth correction, if
$N$ packets arrive per epoch:
\begin{equation}\label{eq:epoch-skew}
  \gamma_3[w_k] = \frac{\gamma_3[j]}{\sqrt{N}},
\end{equation}
which shrinks as $1/\sqrt{N}$ but remains nonzero for finite~$N$.

Through the linear controller, the queue error skewness is
\begin{equation}\label{eq:skew-propagation}
  \kappa_3[e] = \kappa_3[w] \cdot \int_0^\infty h_e(\tau)^3\,\dd\tau,
\end{equation}
where $h_e$ is the closed-loop impulse response from noise to
error.  For a stable linear system with all-positive impulse
response, $\int h_e^3 > 0$, so the output skewness has the same
sign as the input.

\begin{remark}[Linear controllers cannot eliminate skewness]
  Skewness is a third-order cumulant.  A linear filter's transfer
  function controls the power spectrum (second-order), not the
  bispectrum (third-order).  To set $\kappa_3[e] = 0$ from nonzero
  input skewness requires $\int h_e^3 = 0$, which for the class of
  stable minimum-phase responses needed here cannot be achieved.
  Independent control of skewness requires a \emph{nonlinear}
  controller element.
\end{remark}

%% ──────────────────────────────────────────────────────────────────────────────
\subsection{Asymmetric Proportional Controller}%
\label{sec:asym-controller}

The simplest nonlinear element that targets skewness is an
\emph{asymmetric gain}: the controller responds more aggressively
when the queue error is in the dangerous (underrun) direction.

Replace the constant proportional gain~$K$ with a piecewise gain:
\begin{equation}\label{eq:asym-gain}
  K(e) \;=\;
  \begin{cases}
    K_+ & \text{if } e \geq 0 \quad\text{(queue above target)},\\[4pt]
    K_- & \text{if } e < 0   \quad\text{(queue below target)},
  \end{cases}
\end{equation}
with $K_- > K_+$ to respond more aggressively in the underrun
direction.

For analytical tractability, we analyse the first-order
(proportional) controller $u = K(e)\,e$ under the fluid
approximation.  The queue error SDE becomes a
\emph{piecewise-linear} diffusion:
\begin{equation}\label{eq:asym-sde}
  \dd e \;=\; -F_s\,K(e)\,e\;\dd\tau
            \;+\; \sigma_w\;\dd W.
\end{equation}
This is an Ornstein--Uhlenbeck process with rate
$\alpha_+ = F_s\,K_+$ for $e > 0$ and rate
$\alpha_- = F_s\,K_-$ for $e < 0$.

%% ──────────────────────────────────────────────────────────────────────────────
\subsection{Fokker--Planck Equation}%
\label{sec:fokker-planck}

Write $D = \sigma_w^2/2$ for the diffusion coefficient.  The
stationary Fokker--Planck equation for the probability density
$p(e)$ is
\begin{equation}\label{eq:fp}
  \frac{\partial}{\partial e}\bigl[\alpha(e)\,e\,p(e)\bigr]
  + D\,\frac{\partial^2 p}{\partial e^2} = 0,
\end{equation}
where $\alpha(e) = \alpha_\pm$ on each half-line.  On each
half, ~\eqref{eq:fp} is the standard Fokker--Planck equation for an
OU process, whose stationary solution is Gaussian:
\begin{equation}\label{eq:fp-solution}
  p(e) \;=\;
  \begin{cases}
    \displaystyle
    C\,\exp\!\Bigl(-\frac{e^2}{2\sigma_+^2}\Bigr)
      & e \geq 0, \\[10pt]
    \displaystyle
    C\,\exp\!\Bigl(-\frac{e^2}{2\sigma_-^2}\Bigr)
      & e < 0,
  \end{cases}
\end{equation}
with half-variances
\begin{equation}\label{eq:half-var}
  \sigma_+^2 = \frac{D}{\alpha_+}
    = \frac{\sigma_w^2}{2\,F_s\,K_+},
  \qquad
  \sigma_-^2 = \frac{D}{\alpha_-}
    = \frac{\sigma_w^2}{2\,F_s\,K_-}.
\end{equation}

\paragraph{Matching conditions.}
At $e = 0$: (i)~continuity of $p$ is ensured by using the same
constant~$C$ on both halves; (ii)~the probability current
$J = -\alpha\,e\,p - D\,p'$ vanishes identically on each half
(since $p$ satisfies the OU equation), ensuring zero-current
continuity.

\paragraph{Normalisation.}
\begin{equation}\label{eq:split-norm}
  1 = C\!\int_0^\infty\!
      e^{-e^2/(2\sigma_+^2)}\dd e
    + C\!\int_{-\infty}^0\!
      e^{-e^2/(2\sigma_-^2)}\dd e
    = C\,\sqrt{\tfrac{\pi}{2}}\,
      (\sigma_+ + \sigma_-)
\end{equation}
gives
$C = \sqrt{2/\pi}\,/\,(\sigma_+ + \sigma_-)$.

\begin{definition}[Split normal distribution]%
\label{def:split-normal}
  The density~\eqref{eq:fp-solution} is a \emph{split normal}
  (two-piece Gaussian) with scale parameters $\sigma_+$ and
  $\sigma_-$.  When $\sigma_+ = \sigma_-$ it reduces to the
  ordinary Gaussian.
\end{definition}

%% ──────────────────────────────────────────────────────────────────────────────
\subsection{Moments of the Split Normal}%
\label{sec:split-moments}

Let $r = \sigma_-/\sigma_+ = \sqrt{K_+/K_-}$ be the asymmetry
ratio ($r < 1$ when the underrun side is tighter).  All moments
are expressible in closed form.

\paragraph{Mean.}
\begin{equation}\label{eq:split-mean}
  \E[e] = \sqrt{\frac{2}{\pi}}\;(\sigma_+ - \sigma_-)
    = \sqrt{\frac{2}{\pi}}\;\sigma_+\,(1 - r),
\end{equation}
which is \emph{positive} when $K_- > K_+$ ($r < 1$): the
asymmetric controller shifts the mean toward the safe (positive)
side.

\paragraph{Second moment and variance.}
\begin{align}
  \E[e^2]
    &= \frac{\sigma_+^3 + \sigma_-^3}{\sigma_+ + \sigma_-}
    = \sigma_+^2\,\frac{1 + r^3}{1 + r},
    \label{eq:split-e2}\\[6pt]
  \Var[e]
    &= \sigma_+^2\left[
        \frac{1 + r^3}{1 + r}
        - \frac{2}{\pi}(1 - r)^2
      \right].
    \label{eq:split-var}
\end{align}

\paragraph{Skewness.}
For small asymmetry $r = 1 - \epsilon$ with $\epsilon \ll 1$,
the skewness coefficient simplifies to:
\begin{equation}\label{eq:split-skew-approx}
  \gamma_3 \;\approx\;
    \sqrt{\frac{2}{\pi}}\;
    \frac{(4/\pi - 1)\,\epsilon}
         {\bigl(1 - 2/\pi\bigr)^{3/2}}
    \;\approx\; 0.31\,\epsilon.
\end{equation}

\begin{remark}[Interpretation]
  Each unit of gain asymmetry ($\epsilon = 1 - r
  = 1 - \sqrt{K_+/K_-}$) produces $\approx 0.31$ units of
  skewness.  Since $\gamma_3$ of the input noise is typically
  $0.3$--$1.0$, a modest asymmetry of $K_-/K_+ \approx 1.5$
  ($\epsilon \approx 0.18$, $\gamma_3 \approx 0.06$) can
  already partially compensate the intrinsic jitter skewness.
\end{remark}

%% ──────────────────────────────────────────────────────────────────────────────
\subsection{Underrun Probability}%
\label{sec:asym-underrun}

The probability that the queue error exceeds a threshold~$-M$
(i.e.\ the queue drops $M$ samples below target) is the left
tail of the split normal:
\begin{equation}\label{eq:asym-underrun}
  P(e < -M)
    = \frac{\sigma_-}{\sigma_+ + \sigma_-}\;
      \operatorname{erfc}\!\left(
        \frac{M}{\sigma_-\sqrt{2}}
      \right).
\end{equation}

For the symmetric controller ($\sigma_+ = \sigma_- = \sigma$):
\begin{equation}
  P_{\text{sym}}(e < -M) = \tfrac{1}{2}\,
    \operatorname{erfc}\!\bigl(M/(\sigma\sqrt{2})\bigr).
\end{equation}

The key observation is that in the asymmetric case, the argument
of the erfc is $M/\sigma_-$, and since $\sigma_- < \sigma_+$ (the
underrun side is tighter), the effective number of standard
deviations \emph{increases}:
\begin{equation}
  \frac{M}{\sigma_-}
    = \frac{M}{\sigma_+\,r}
    = \frac{1}{r}\,\frac{M}{\sigma_+}.
\end{equation}

\begin{proposition}[Underrun reduction]%
\label{prop:underrun-reduction}
  With asymmetry ratio $r = \sqrt{K_+/K_-}$, a buffer margin
  of~$M$ provides $M/(\sigma_+\,r)$ effective standard deviations
  of protection on the underrun side, compared to $M/\sigma$ for
  the symmetric case.  For $K_-/K_+ = 2$ ($r = 1/\sqrt{2}$):
  \begin{itemize}
    \item A $3\sigma$ margin becomes effectively $4.2\sigma$,
    \item The underrun probability drops by a factor of~${\approx\,10}$.
  \end{itemize}
\end{proposition}

%% ──────────────────────────────────────────────────────────────────────────────
\subsection{The Variance--Skewness Tradeoff}%
\label{sec:var-skew-tradeoff}

The asymmetric controller shifts the mean to $\E[e] > 0$
(the queue sits slightly above target on average).  This is a
deliberate asymmetric operating point that protects the underrun
tail.

For small asymmetry, the variance penalty is \emph{second-order}
while the underrun benefit is \emph{first-order}:

\begin{proposition}[Variance--skewness tradeoff]%
\label{prop:var-skew}
  At leading order in $\epsilon = 1 - r$:
  \begin{itemize}
    \item Underrun probability decreases exponentially
      (via the increased $M/\sigma_-$ ratio),
    \item Variance increases as $1 + O(\epsilon^2)$
      (quadratic cost),
    \item Mean shift is $O(\epsilon)$ (linear, compensated by
      integral action if present).
  \end{itemize}
  The tradeoff strongly favours small asymmetry when the
  buffer margin~$M$ is large relative to~$\sigma$.
\end{proposition}

%% ──────────────────────────────────────────────────────────────────────────────
\subsection{Extension to the PID Controller}%
\label{sec:asym-extension}

The Fokker--Planck analysis above uses a first-order
(proportional) controller for analytical tractability.  For the
PID controller (Section~\ref{sec:bias}), the state space is
three-dimensional $(e, \xi, u)$ and the Fokker--Planck equation
becomes a 3D PDE, which does not admit closed-form solutions in
general.

However, the qualitative results carry over:
\begin{enumerate}
  \item The marginal density of~$e$ remains approximately
    split-normal when the gain asymmetry is a slowly-varying
    modulation of the linear dynamics (\emph{adiabatic
    approximation}).
  \item The integral action ($K_0$) compensates the mean shift
    introduced by the asymmetric gain, restoring $\E[e] = 0$
    while preserving the asymmetric tail structure.
  \item The dominant pole separation ($p_0 \ll \omega_n$) means
    the integral mode operates on a much slower timescale
    and does not interact with the asymmetric gain dynamics.
\end{enumerate}

In practice, the asymmetric PID controller would use:
\begin{equation}\label{eq:asym-pid}
  v \;=\;
    K_1(e)\,e + K_0\,\xi - K_2\,u,
  \qquad
  K_1(e) = K_1\,(1 + \beta\,\operatorname{sgn}(-e)),
\end{equation}
where $\beta \in [0, 1)$ controls the asymmetry.
The effective gains are $K_- = K_1(1+\beta)$ for $e < 0$
and $K_+ = K_1(1-\beta)$ for $e > 0$, giving
$r = \sqrt{(1-\beta)/(1+\beta)}$.

%% ══════════════════════════════════════════════════════════════════════════════
%\bibliographystyle{plain}
%\bibliography{references}

\end{document}
